\begin{figure}[H]
\centering
\begin{tikzpicture}[>=Latex,scale=1.0,font=\small]
\draw[blue!60!black,thick] (0,0) circle (1.9);
\draw[gray!55,dashed] (1.9,-2.1)--(1.9,2.4);
\fill (0,0) circle(1.5pt) node[below left] {\(0\)};
\coordinate (tag) at (1.9,0);
\coordinate (candidate) at (1.9,1.65);
\draw[->,thick] (0,0)--(tag) node[midway,below] {\(t_s\)};
\draw[->,thick,orange!75!black] (tag)--(candidate) node[midway,right] {\(d=u-u'\)};
\draw[->,thick,gray!75!black] (0,0)--(candidate) node[midway,above left] {\(c=t_s+d\)};
\fill (tag) circle(2pt) node[below right] {\(t_s\)};
\fill (candidate) circle(2pt) node[above right] {\(c\)};
\draw (1.65,0)--(1.65,.25)--(1.9,.25);
\node[align=left,anchor=west,text width=6cm] at (4.1,0.6) {
The affine slice gives \(t_s\cdot d=0\).\\[5pt]
The group equation gives \(c=t_s+d\).\\[5pt]
If \(c\) is another tag on the same sphere,\\[3pt]
\(\|c\|^2=\|t_s\|^2+\|d\|^2=\|t_s\|^2\),\\[3pt]
so \(d=0\) and \(c=t_s\).
};
\end{tikzpicture}
\caption{The separation identity in a two-dimensional slice.
The dashed line represents the displacement hyperplane
\(t_s+t_s^\perp\), rather than the slice containing \(u,u'\).
A nonzero displacement along it leaves the tag sphere, so no other tag
can survive. The argument applies in every dimension.}
\label{fig:group-geometry}
\end{figure}
